
\subsubsection*{2.4.A}
\phantomsection
\label{sga3:VIIB:2.4.A}

设 \(G\) 和 \(Q\) 为拓扑平坦的形式 \(k\)-群。假设存在同态 \[
  \sigma:Q\longrightarrow G
  \qquad\text{且}\qquad
  \pi:G\longrightarrow Q
\] 使得 \[
  \pi\circ\sigma=\operatorname{id}_Q.
\]。\footnote{\label{sga3:VIIB:nde:92}编者添加了段落 2.4.A 和 \SGARefLink{sga3:VIIB:2.4.B}{2.4.B}，它们分别将用于 \SGARefLink{sga3:VIIB:5.1.3}{5.1.3} 和 \SGARefLink{sga3:VIIB:2.9}{2.9}。}特别地，\(\sigma\) 为单态射，故 \(Q\) 为 \(G\) 的形式子群 \textup{（参见 \SGARefLink{sga3:VIIB:1.3}{注 1.3}）}。设 \[
  N=\operatorname{Ker}(\pi)
\]，且 \(\sigma'\) 为包含映射 \(N\hookrightarrow G\)。则 \(G\) 为 \(N\) 被 \(Q\) 扩张的半直积 \textup{（参见 I，\SGARefLink{sga3:i:proposition:2-3-8}{2.3.8}）}。换言之，对每个 \[
  B\in\operatorname{Ob}(\operatorname{Alf}/k),
\]，通过 \(\sigma'\) 和 \(\sigma\) 将 \(N(B)\) 和 \(Q(B)\) 与其在 \(G(B)\) 中的像等同。则 \(N(B)\) 为 \(G(B)\) 的正规子群，且映射 \begin{equation}
  \mu:N(B)\times Q(B)\longrightarrow G(B),
  \qquad
  (x,q)\longmapsto xq
  \tag{1}
  \label{sga3:VIIB:2.4.A:eq1}
\end{equation} 为双射。

于是形式 \(k\)-簇态射 \begin{equation}
  \theta:G(B)\longrightarrow N(B),
  \qquad
  g\longmapsto g\cdot\sigma\pi(g^{-1})
  \tag{2}
  \label{sga3:VIIB:2.4.A:eq2}
\end{equation} 为 \(\sigma'\) 的收缩，且 \(\mu\) 的逆为 \begin{equation}
  g\longmapsto\bigl(\theta(g),\pi(g)\bigr).
  \tag{3}
  \label{sga3:VIIB:2.4.A:eq3}
\end{equation}。若 \(p:G\to G/Q\) 表示商投影，则 \[
  p\circ\sigma':N\longrightarrow G/Q
\] 为形式 \(k\)-簇的同构。

\smallskip\noindent\emph{原文说明。}Polo--Gille 文本此处印作 \(\theta\circ\sigma':N\to G/Q\)。由于 \(\theta:G\to N\)，该复合的类型不匹配；具有所显示源与目标的同构应为商复合 \(p\circ\sigma'\)。

\smallskip 特别地，由 \SGARefLink{sga3:VIIB:1.4.thm}{1.4} \textup{\SGARefLink{sga3:vii:paragraph:5-1:equation:ii}{(ii)}}，\(N\) 在 \(k\) 上拓扑平坦。设 \[
  \alpha:Q\times N\longrightarrow N
\] 和 \[
  \beta:G\times N=N\times Q\times N\longrightarrow N
\] 为在点上由 \[
  \alpha(q,y)=qyq^{-1},
  \qquad
  \beta(x,q,y)=x\alpha(q,y).
\] 定义的映射。则下图交换：\begin{center}\phantomsection \label{sga3:VIIB:2.4.A:eq4-diagram} % SGA3-VII-DIAG-109
% Authority: Exp7B-13oct24.pdf, local p. 53.
% Crop: figures/Exp7B_localp053_semidirect_product_diagram4.png
% Locator: tex/components/43_expose_VIIB_section24A_semidirect_products_en.tex, Loop 1 line 101.
\begin{tikzcd}[row sep=3.2em,column sep=3.5em]
{} \arrow[d,phantom,"\textup{(4)}" description]
  & G\times G
  \arrow[r,"{\operatorname{id}\times\theta}"]
  \arrow[d,"m_G"']
  & G\times N
      \arrow[r,"{\theta\times\operatorname{id}}"]
      \arrow[d,"\beta"]
  & N\times N \arrow[dl,"m_N"] \\
{} & G \arrow[r,"\theta"] & N & {.}
\end{tikzcd}
\end{center}

用 \(G\)、\(Q\) 和 \(N\) 的仿射代数 \(A\)、\(A_0\) 和 \(A'\)，可以将这一点表述如下 \textup{（参见下面的 \SGARefLink{sga3:VIIB:5.1.3}{5.1.3}）}。设 \[
  \rho':A\longrightarrow A',
  \qquad
  \rho:A\longrightarrow A_0,
  \qquad
  \tau:A_0\longrightarrow A
\] 为对应于 \(\sigma'\)、\(\sigma\) 和 \(\pi\) 的 \(k\)-双代数同态，并置 \[
  I=\operatorname{Ker}(\rho).
\]。由前面的推论，\(A'\) 与 \[
  A\big/\overline{A\,\tau(J_0)},
\] 等同，其中 \(J_0\) 是 \(A_0\) 的增广理想。

另一方面，设 \(B\) 为 \(G/Q\) 的仿射代数。它是态射对 \phantomsection \label{sga3:VIIB:2.4.A:affine-equalizer-pair} \begin{equation*}
  A
  \mathrel{\overset{\tau_1}
    {\underset{(\operatorname{id}_A\widehat\otimes\rho)
      \circ\Delta_A}{\rightrightarrows}}}
  A\widehat\otimes_kA_0,
\end{equation*} 的核，其中 \(\tau_1(x)=x\widehat\otimes1\)。故 \[
  B
  =
  \bigl\{
    x\in A
    \mathrel{\big|}
    \Delta_A(x)-x\widehat\otimes1
    \in A\widehat\otimes I
  \bigr\}.
\]。设 \[
  \gamma=\theta^\natural:A'\longrightarrow A
\] 为相应的连续 \(k\)-代数同态。它是 \(\rho'\) 的截面，且为从 \(A'\) 到 \(B\) 的预有限 \(k\)-代数同构。此外，\[
  \tau=\pi^\natural
\] 将 \(A_0\) 等同于 \(A\) 的子双代数，即商 \(N\backslash G\) 的仿射代数。

由式~\eqref{sga3:VIIB:2.4.A:eq1} 和~\eqref{sga3:VIIB:2.4.A:eq3} 可得预有限 \(k\)-代数同构 \begin{equation}
  A'\widehat\otimes_kA_0
  \xrightarrow{\ \sim\ }
  A,
  \qquad
  a'\widehat\otimes a_0
  \longmapsto
  \gamma(a')\tau(a_0),
  \tag{*}
  \label{sga3:VIIB:2.4.A:star}
\end{equation}，其逆为映射 \[
  a\longmapsto
  (\rho'\widehat\otimes\rho)\Delta_A(a).
\]。

最后，通过 \(\gamma\) 将 \(A'\) 与其在 \(A\) 中的像等同，于是投影 \(A\to A'\) 就是 \(\gamma\rho'\)。以 \(\Delta_{A'}\) 表示 \(A'\) 的余乘法，由 \hyperref[sga3:VIIB:2.4.A:eq4-diagram]{\textup{(4)}} 可得 \[
  \Delta_A(A')\subset A\widehat\otimes A'
\]，且下图交换：\begin{center}\phantomsection \label{sga3:VIIB:2.4.A:Aprime-comultiplication-square} % SGA3-VII-DIAG-110
% Authority: Exp7B-13oct24.pdf, local p. 53.
% Crop: figures/Exp7B_localp053_Aprime_comultiplication_square.png
% Locator: tex/components/43_expose_VIIB_section24A_semidirect_products_en.tex, Loop 1 line 193.
\begin{tikzcd}[row sep=3em,column sep=3.7em]
A' \arrow[r,"\Delta_A"] \arrow[d,equals]
  & A\widehat\otimes A' \arrow[d,"{\gamma\rho'\widehat\otimes
                                      \operatorname{id}}" right] \\
A' \arrow[r,"\Delta_{A'}"]
  & A'\widehat\otimes A'
\end{tikzcd}
\end{center}因此还有 \[
  \gamma\rho'\circ c_A=c_{A'},
\]，其中 \(c_A\) 和 \(c_{A'}\) 分别为 \(A\) 和 \(A'\) 的对极。另一方面，以 \(m_A\) 表示 \(A\) 中的乘法，式~\eqref{sga3:VIIB:2.4.A:eq2} 对每个 \(a\in A\) 给出 \[
  \gamma\rho'(a)
  =
  \bigl(
    m_A
    \circ
    (\operatorname{id}_A\widehat\otimes\tau\rho c_A)
    \circ
    \Delta_A
  \bigr)(a).
\]。
